\subsection{Case study: the QCD axion and neutrino masses}
\label{case:2608-05150}

\paragraph{Paper identity and theme.}
The source is Hridoy Debnath, Pavel Fileviez P\'erez, and Grace Miller,
``The QCD Axion and Neutrino Masses,'' arXiv:2608.05150v1. The paper links
Peccei--Quinn (PQ) breaking to neutrino-mass generation in DFSZ and KSVZ
models. It compares anomaly coefficients, derives axion couplings and decay
rates, and studies neutrino signals from axion dark-matter decay. Statements
in the next two paragraphs report what the paper says; paragraphs explicitly
headed ``Interface analysis'' report the present stress-test analysis.

\paragraph{Source census method.}
The census used the official v1 archive identified by the frozen manifest and
read \texttt{tmp/claim-interface-study/extracted/2608.05150/main.tex} from its
macro definitions through both occurrences of \verb|\end{document}|. A
claim-bearing source unit is a prose paragraph, list item, equation-plus-lead,
table, or figure caption that asserts a definition, assumption, imported
result, derivation step, result, limitation, or interpretation. Pure headings,
layout code, acknowledgements, and bibliographic commands were excluded.
Repeated abstract, introduction, body, and summary statements were counted as
occurrences and then consolidated. Material after the first
\verb|\end{document}| was retained as ``dormant source'' rather than silently
presented as compiled appendix content.

\begin{center}
\begin{tabular}{lrr}
\hline
Source division & Claim-bearing units & Consolidated candidates \\
\hline
Abstract & 3 & 3 \\
Introduction & 8 & 7 \\
The QCD axion & 6 & 3 \\
DFSZ mechanisms & 17 & 9 \\
KSVZ mechanisms and colored seesaw & 16 & 8 \\
Neutrinos from the QCD axion & 18 & 11 \\
Summary & 5 & 0 (repeated support) \\
Dormant post-document derivations & 10 & 3 \\
\hline
Total occurrences & 83 & 35 unique candidates \\
\hline
\end{tabular}
\end{center}

The candidate count is not the sum of the right column because several
candidates join non-contiguous support across divisions, while several source
units split into more than one candidate. The JSONL inventory is the
candidate-level source of truth.

\paragraph{Compact claim-family census.}
\begin{center}
\begin{tabular}{p{0.24\linewidth}p{0.48\linewidth}r}
\hline
Family & Examples retained & Candidates \\
\hline
Background and imported results & Strong CP bound, neutrino-mass evidence,
seesaw and radiative mechanisms & 3 \\
Model definitions and assumptions & DFSZ and KSVZ field content, shared-scale
hypothesis, model scope & 6 \\
Analytical relations & $m_a(f_a)$, $g_{a\gamma\gamma}$, seesaw and loop masses,
anomaly maps, decay widths, flux & 13 \\
Numerical and phenomenological findings & Current exclusions, branching and
lifetime patterns, spectra, flux comparisons & 8 \\
Interpretation, testability, and scope & Dark-matter signal and PTOLEMY-type
testability & 2 \\
Dormant derivations & DFSZ color and electromagnetic anomaly calculations & 3 \\
\hline
\end{tabular}
\end{center}

\paragraph{Representative difficult splits.}
The paper places a model-family assertion, a coefficient convention, and a
mass scaling in nearby prose. The inventory separates the explicit relations
from the family-level assertion:
\[
\begin{aligned}
 m_a &\simeq 5.7\,\mu\mathrm{eV}
 \left(\frac{10^{12} \mathrm{GeV}}{f_a}\right),\\
 g_{a\gamma\gamma} &= \frac{\alpha_{\mathrm{em}}}{2\pi f_a}
 \left(\frac{E}{N}-1.92\right).
\end{aligned}
\]
The first two normalize directly. The separate claim that all linked models
have axion--neutrino interactions is retained with its universal quantifier and
model scope rather than being inferred from these equations.

A second split concerns axion decay. The two channel formulas are one
analytical candidate,
\[
\begin{aligned}
 \Gamma(a\to\gamma\gamma)
 &=\frac{\alpha_{\mathrm{em}}^2}{256\pi^3}
 \left(\frac EN-1.92\right)^2\frac{m_a^3}{f_a^2},\\
 \Gamma(a\to\nu_i\nu_i)
 &=\frac{m_a}{16\pi}
 \left|\frac{c_\nu m_{\nu_i}}{Nf_a}\right|^2
 \sqrt{1-\frac{4m_{\nu_i}^2}{m_a^2}},
\end{aligned}
\]
whereas channel dominance for normal or inverted ordering is a separate,
figure-supported candidate with the massless-lightest-neutrino assumption.
The statement that the axion remains a good dark-matter candidate is yet
another interpretive conclusion; it is not folded into either width formula.

The flux claim also needs an assumption-bearing split:
\[
\begin{aligned}
 \frac{d\Phi_\nu}{dE}
 &=\frac{1}{4\pi m_a\tau_a}\frac13\frac{dN}{dE}D(\Omega),\\
 \frac{dN}{dE}&=2\delta(E-m_a/2).
\end{aligned}
\]
This formula is normalized conditionally on axions constituting Galactic dark
matter. The claims that the flux is ``large'', lies above backgrounds, or is
testable are retained separately because they require numerical curves,
background definitions, and detector assumptions.

\paragraph{Source-fidelity decisions.}
Every JSONL span uses one-indexed complete source lines and byte-for-byte text
from those lines. Repeated conclusions were consolidated with multiple spans
when the later occurrence supplied truth-relevant scope. Knowledge
attribution distinguishes community background, cited imported formulas, and
current-work model or scan results. Approximate signs, conditional phrases
such as ``assuming'', and feasibility language such as ``could in principle''
were preserved. No candidate infers experimental reach beyond the paper's
wording. The post-document anomaly calculations are explicitly designated as
inactive source content: their presence is evidence that source state cannot
be represented by line location alone.

\paragraph{Interface analysis: MathClaimIR boundaries.}
The 35 candidates divide into 16 \textsc{yes}, 14 \textsc{partial}, and 5
\textsc{no} decisions. Explicit equations, inequalities, decay widths, and
finite model-to-anomaly tables are \textsc{yes}. Family-wide anomaly claims,
shared-scale explanations, and hierarchy-dependent channel statements are
\textsc{partial}: their mathematical core is useful, but model membership,
assumptions, and causal explanation remain semantic. Figure-level exclusions,
``clearly separated'' fluxes, dark-matter viability, and experimental
feasibility are \textsc{no}; forcing these into a formula would erase the
comparison metric or evidence provenance.

\paragraph{Interface pressures.}
This paper stresses the interface in six recurring ways. First, one atomic
claim may need non-contiguous support. Second, a universal model-family claim
must not inherit stronger scope from one worked example. Third, table cells
are evidence, not merely presentation. Fourth, figure-dependent comparisons
need links to external contours and numerical assumptions. Fifth, anomaly
coefficients depend on charge and normalization conventions. Sixth, dormant
TeX after \verb|\end{document}| requires source-activity status distinct from
scientific acceptance. These are pressures on provenance, scope, and evidence
completeness; they do not justify enlarging the proposed Registry schema in
this case study.

\paragraph{Provisional verdict.}
The interface handles the paper's equation-level content cleanly when an
atomic statement carries explicit scope and exact source spans. It is only
partially adequate for model-family conclusions and numerical comparisons,
and it must allow \textsc{no} normalization without treating the scientific
claim as invalid. The decisive requirement is separation of source fidelity,
mathematical normalization, and scientific or phenomenological acceptance.